% Copyright 2026 Torbjörn E. M. Nordling
% SPDX-License-Identifier: Apache-2.0
% frame-pipeline.tex — From raw video to figures (18:30–20:00)
% Data README "Pipeline": track, process, check, analyse; Sam's guide §7-8.
% Conceptual only: no script names (course rule).
\begin{frame}{From video to figure}{Five steps; the video itself is never stored in git}
  \footnotesize
  \centering
  \begin{tikzpicture}[node distance=5mm]
    \node[bd] (rec) {Record\\video + Pi log};
    \node[bd, right=of rec] (trk) {Track\\video $\to$ frame table};
    \node[bd, right=of trk] (prc) {Process\\angle, speed, duty, $\sigma$};
    \node[bd, right=of prc] (chk) {Check\\quality report};
    \node[bd, right=of chk] (ana) {Analyse\\fits, figures, tables};
    \draw[bdline] (rec) -- (trk);
    \draw[bdline] (trk) -- (prc);
    \draw[bdline] (prc) -- (chk);
    \draw[bdline] (chk) -- (ana);
  \end{tikzpicture}

  \vspace{6pt}
  \raggedright
  \begin{itemize}\setlength{\itemsep}{2pt}
    \item \textbf{Record:} copy the video by cable or card, never via a messaging app (it compresses the video and changes its timing). Copy the Pi logs unchanged and compare checksums.
    \item \textbf{Track:} reduces each 1--2\,GB video to a table of marker, hub, and LED pixels per frame, about 1\,MB. This table is your \alert{raw data}.
    \item \textbf{Process:} computes angle, speed, and applied duty with uncertainty, and aligns the clocks by the LED.
    \item \textbf{Check:} read the quality report. A warning means: use another trial of the same condition, or record the clip again. Trials must be synchronised by the LED.
    \item \textbf{Analyse:} fits and figures are \alert{regenerated by code}, never edited by hand.
  \end{itemize}
\end{frame}
